\documentclass[a4paper,11pt]{article}

\usepackage{amsmath}
\usepackage{amssymb}
\usepackage{mathtools}
\usepackage{tabularx}
\usepackage{booktabs}
\usepackage{mdframed}
\usepackage{microtype}
\usepackage{hyperref}
\usepackage{xcolor}
\usepackage{geometry}

\geometry{margin=1in}

% Number sets
\newcommand*{\N}{\mathbb N}
\newcommand*{\Z}{\mathbb Z}
\newcommand*{\Q}{\mathbb Q}
\newcommand*{\R}{\mathbb R}

% Math operators
\DeclareMathOperator*{\dom}{dom}
\DeclareMathOperator*{\argmin}{arg\,min}
\DeclareMathOperator*{\argmax}{arg\,max}

% Delimiters
\DeclarePairedDelimiter\abs\lvert\rvert
\DeclarePairedDelimiter\seq\langle\rangle
\providecommand\given{}
\newcommand\SetSymbol[1][]{%
	\nonscript\:#1\vert
	\allowbreak
	\nonscript\:
	\mathopen{}}
\DeclarePairedDelimiterX\Set[1]\lbrace\rbrace{%
	\renewcommand\given{\SetSymbol[\delimsize]}
	#1
}

% Counters for framed environments, reset per section
\newcounter{defcount}[section]
\renewcommand{\thedefcount}{\thesection.\arabic{defcount}}
\newcounter{gramcount}[section]
\renewcommand{\thegramcount}{\thesection.\arabic{gramcount}}

% Framed environments
\mdfdefinestyle{definitionstyle}{%
  linecolor=black,linewidth=1pt,%
  backgroundcolor=yellow!10,
  innertopmargin=6pt,
  innerbottommargin=8pt,
  roundcorner=5pt,
}

\mdfdefinestyle{grammarstyle}{%
  linecolor=black,linewidth=1pt,%
  backgroundcolor=lightgray!5,
  innertopmargin=6pt,
  innerbottommargin=8pt,
  roundcorner=5pt,
  font=\normalfont,
}

\newenvironment{frameddef}[1][]{%
  \stepcounter{defcount}%
  \begin{mdframed}[style=definitionstyle]%
  {\normalfont\bfseries Definition~\thedefcount\ifx\relax#1\relax\else\ (#1)\fi.}\par\vspace{4pt}\noindent\ignorespaces
}{%
  \end{mdframed}%
}

\newenvironment{framedgram}[1][]{%
  \stepcounter{gramcount}%
  \begin{mdframed}[style=grammarstyle]%
  {\normalfont\bfseries Grammar~\thegramcount\ifx\relax#1\relax\else\ (#1)\fi.}\par\vspace{4pt}\noindent\ignorespaces
}{%
  \end{mdframed}%
}

% OPQL syntax highlighting commands
\newcommand{\opqltoken}[1]{\textcolor{teal}{#1}}
\newcommand{\opqlevent}[1]{\textcolor{orange}{E(}\textcolor{black}{#1}\textcolor{orange}{)}}
\newcommand{\opqlobject}[1]{\textcolor{cyan}{O(}\textcolor{black}{#1}\textcolor{cyan}{)}}
\newcommand{\opqlrela}[1]{\textcolor{gray}{-[}\textcolor{black}{#1}\textcolor{gray}{]-}}
\newcommand{\opqlrelr}[1]{\textcolor{gray}{-[}\textcolor{black}{#1}\textcolor{gray}{]->}}
\newcommand{\opqlrell}[1]{\textcolor{gray}{<-[}\textcolor{black}{#1}\textcolor{gray}{]-}}

\newcommand{\nterm}[1]{\textit{#1}}
\newcommand{\term}[1]{\texttt{\color{blue}#1}}

% Title
\title{OPQL Language Standard}
\author{}
\date{}

\begin{document}

\maketitle

\tableofcontents
\newpage

\section{Introduction}

OPQL (Object-centric Process Query Language) is a query language for filtering object-centric event logs conforming to the OCEL 2.0 standard and extracting tabular data from them. This document defines the complete syntax and semantics of OPQL version 1.0. The language supports pattern matching over events, objects, and their qualified relations; arithmetic and logical expressions; temporal reasoning; and aggregation via subqueries. Queries are evaluated sequentially against an event log, transforming an initial empty context into a result that is either a filtered event log or a table of records.

\section{Preliminaries}

This section establishes the mathematical foundations and the OCEL 2.0 metamodel upon which OPQL operates.

\subsection{Sets and Multisets}

A \emph{set} is a collection of distinct elements. We denote sets using curly braces $\{ \dots \}$. The \emph{powerset} of a set $A$, denoted $\mathbb{P}(A)$, is the set of all subsets of $A$.

A \emph{multiset} (or bag) over a set $M$ is a tuple $\mathcal{M} = (M, f)$ where $f : M \to \mathbb{N}$ is a multiplicity function. We denote multisets using square brackets $[ \dots ]$. The \emph{bag union} of multisets $\mathcal{M} = (M, f)$ and $\mathcal{N} = (M, g)$ is $\mathcal{M} \uplus \mathcal{N} = (M, h)$ with $h(x) = f(x) + g(x)$.

\subsection{Partial Functions}

A \emph{partial function} $f : X \nrightarrow Y$ is a function from a subset $S \subseteq X$ to $Y$. For $a \in X \setminus S$, we write $f(a) = \bot$. In OPQL, $\bot$ is denoted as $None$. Partial functions arise wherever a lookup may have no result: an event may not have a particular attribute, an object attribute may not be defined at a given timestamp, and so on.

\subsection{OCEL 2.0 Universes}

Let $\mathbb{U}_\Sigma$ be the universe of strings. The following pairwise disjoint universes are defined:

\begin{itemize}
    \item $\mathbb{U}_{ev} \subseteq \mathbb{U}_\Sigma$ --- universe of event identifiers
    \item $\mathbb{U}_{etype} \subseteq \mathbb{U}_\Sigma$ --- universe of event types
    \item $\mathbb{U}_{obj} \subseteq \mathbb{U}_\Sigma$ --- universe of object identifiers
    \item $\mathbb{U}_{otype} \subseteq \mathbb{U}_\Sigma$ --- universe of object types
    \item $\mathbb{U}_{attr} \subseteq \mathbb{U}_\Sigma$ --- universe of attribute names
    \item $\mathbb{U}_{val}$ --- universe of attribute values
    \item $\mathbb{U}_{time}$ --- universe of timestamps (with $0 \in \mathbb{U}_{time}$ as the smallest element and $\infty \in \mathbb{U}_{time}$ as the largest)
    \item $\mathbb{U}_{qual} \subseteq \mathbb{U}_\Sigma$ --- universe of relation qualifiers
\end{itemize}

\subsection{Object-Centric Event Log}

An Object-Centric Event Log (OCEL) is a tuple
\[
L = (E, O, EA, OA, evtype, time, objtype, eatype, oatype, eaval, oaval, E2O, O2O)
\]
with:
\begin{itemize}
    \item $E \subseteq \mathbb{U}_{ev}$ --- set of events
    \item $O \subseteq \mathbb{U}_{obj}$ --- set of objects
    \item $EA \subseteq \mathbb{U}_{attr}$ --- set of event attributes
    \item $OA \subseteq \mathbb{U}_{attr}$ --- set of object attributes
    \item $evtype : E \to \mathbb{U}_{etype}$ --- assigns types to events
    \item $time : E \to \mathbb{U}_{time}$ --- assigns timestamps to events
    \item $objtype : O \to \mathbb{U}_{otype}$ --- assigns types to objects
    \item $eatype : EA \to \mathbb{U}_{etype}$ --- assigns event types to event attributes
    \item $oatype : OA \to \mathbb{U}_{otype}$ --- assigns object types to object attributes
    \item $eaval : (E \times EA) \nrightarrow \mathbb{U}_{val}$ --- assigns values to event attributes
    \item $oaval : (O \times OA \times \mathbb{U}_{time}) \nrightarrow \mathbb{U}_{val}$ --- assigns values to object attributes at a given timestamp
    \item $E2O \subseteq E \times \mathbb{U}_{qual} \times O$ --- qualified event-to-object relations
    \item $O2O \subseteq O \times \mathbb{U}_{qual} \times O$ --- qualified object-to-object relations
\end{itemize}

such that:
\begin{itemize}
    \item $\dom(eaval) \subseteq \{ (e, ea) \in E \times EA \mid evtype(e) = eatype(ea) \}$
    \item $\dom(oaval) \subseteq \{ (o, oa, t) \in O \times OA \times \mathbb{U}_{time} \mid objtype(o) = oatype(oa) \}$
\end{itemize}

The domain constraints ensure that event attributes are only defined for events of the matching type, and object attributes are only defined for objects of the matching type. Object attributes are versioned: $oaval(o, oa, t)$ gives the value of attribute $oa$ for object $o$ at time $t$; a new entry appears whenever the attribute changes. Relations in $E2O$ associate events with the objects they involve; the qualifier in the triple carries the relationship type (e.g., role of the object in the event). Relations in $O2O$ express direct connections between objects, similarly qualified.

\subsection{Durations}

A duration is the signed difference between two points in time:
\[
\mathbb{U}_{duration} = \{ t_1 - t_2 \mid t_1, t_2 \in \mathbb{U}_{time} \}.
\]
Durations support addition, subtraction, and scaling (multiplication and division by real numbers). A timestamp may be offset by adding or subtracting a duration, yielding a new timestamp. Dividing two durations yields a real number.

\section{Lexical Structure}

This section defines the lexical tokens of OPQL. Whitespace and comments are recognized but discarded during lexical analysis.

\subsection{Keywords}

The following reserved keywords are recognized case-sensitively:

\begin{center}
\begin{tabular}{llll}
\toprule
Keyword & Alternatives & Keyword & Alternatives \\
\midrule
\term{PATTERN} & & \term{FILTER} & \\
\term{SUBJECTTO} & \term{ST} & \term{RETURN} & \\
\term{OCEL} & & \term{KEEP} & \\
\term{DISTINCT} & & \term{ORDERBY} & \\
\term{ASC} & \term{ASCENDING} & \term{DESC} & \term{DESCENDING} \\
\term{LIMIT} & & \term{BINNED} & \\
\term{WHEN} & & \term{AS} & \\
\term{MATERIALIZED} & & & \\
\term{True} & & \term{False} & \\
\term{None} & & \term{E} & \\
\term{O} & & \term{T} & \\
\term{D} & & \term{OR} & \\
\term{XOR} & & \term{AND} & \\
\term{NOT} & & & \\
\bottomrule
\end{tabular}
\end{center}

\subsection{Symbolic Names}

A symbolic name is an identifier used to bind values. It must start with a letter (uppercase or lowercase), followed by any sequence of letters, digits, or underscores.

\begin{framedgram}[Symbolic Names]
\begin{tabularx}{\linewidth}{l c X}
\nterm{symbolicName} & = & (\term{a}..\term{z} | \term{A}..\term{Z}) \{ \term{a}..\term{z} | \term{A}..\term{Z} | \term{0}..\term{9} | \term{\_} \} \\
\end{tabularx}
\end{framedgram}

The set of all valid symbolic names is denoted $\mathbb{U}_{sym}$.

\subsection{Strings}

A string literal is a sequence of characters enclosed in double quotation marks (\term{"}). The following escape sequences are recognized:

\begin{center}
\begin{tabular}{ll}
\toprule
Escape & Meaning \\
\midrule
\term{\textbackslash"} & Double quotation mark \\
\term{\textbackslash\textbackslash} & Backslash \\
\term{\textbackslash/} & Forward slash \\
\term{\textbackslash b} & Backspace \\
\term{\textbackslash f} & Form feed \\
\term{\textbackslash n} & Newline \\
\term{\textbackslash r} & Carriage return \\
\term{\textbackslash t} & Tab \\
\bottomrule
\end{tabular}
\end{center}

\subsection{Punctuation and Operators}

\begin{center}
\begin{tabular}{lll}
\toprule
Token & Name & Usage \\
\midrule
\term{,} & Comma & Separator \\
\term{:} & Colon & Type specifier \\
\term{(} \term{)} & Parentheses & Grouping, function calls \\
\term{[} \term{]} & Square brackets & Property lookup, relations \\
\term{*} & Asterisk & Wildcard, multiplication \\
\term{@} & At-sign & Timestamp qualifier for properties \\
\term{==} & Equal to & Comparison \\
\term{!=} & Not equal & Comparison \\
\term{<=} & Less than or equal & Comparison \\
\term{>=} & Greater than or equal & Comparison \\
\term{<} & Less than & Comparison, left-directed relation \\
\term{>} & Greater than & Comparison, right-directed relation \\
\term{+} & Plus & Addition, unary positive \\
\term{-} & Minus & Subtraction, unary negative, relation anchor \\
\term{/} & Slash & Division \\
\term{\%} & Percent & Modulo \\
\term{\textasciicircum} & Caret & Exponentiation \\
\term{->} & Right arrow & Right-directed relation \\
\term{<-} & Left arrow & Left-directed relation \\
\bottomrule
\end{tabular}
\end{center}

\subsection{Whitespace and Comments}

Whitespace (spaces, tabs, line breaks, and Unicode whitespace characters) is ignored except where needed to separate tokens.

Comments are supported in two forms:
\begin{itemize}
    \item Line comments: \term{//} followed by any characters until end of line
    \item Block comments: \term{/*} followed by any characters until \term{*/}
\end{itemize}

\section{Data Types}

The universe of values OPQL operates on is:
\[
\mathbb{U}_V = \mathbb{U}_\Sigma \cup \mathbb{R} \cup \mathbb{U}_{time} \cup \mathbb{U}_{duration} \cup \mathbb{U}_{bool} \cup \{ None \}
\]
where $\mathbb{U}_\Sigma$ is the universe of strings, $\mathbb{R}$ the real numbers, $\mathbb{U}_{time}$ the universe of timestamps, $\mathbb{U}_{duration}$ the universe of durations, $\mathbb{U}_{bool} = \{ True, False \}$ the universe of booleans, and $None$ the non-value.

\subsection{Strings}

Strings are sequences of encodable characters enclosed in double quotation marks. The characters \term{"} and \term{\textbackslash} must be escaped as \term{\textbackslash"} and \term{\textbackslash\textbackslash} respectively.

\subsection{Integers and Floats}

\begin{framedgram}[Numeric Literals]
\begin{tabularx}{\linewidth}{l c X}
\nterm{int} & = & \term{0} | (\term{1}..\term{9}) \{ \term{0}..\term{9} \} \\
\nterm{float} & = & \nterm{int} \term{.} (\term{0}..\term{9}) \{ \term{0}..\term{9} \} \\
\end{tabularx}
\end{framedgram}

Integers are either zero or a non-zero digit followed by any number of digits. Leading zeros are not permitted. Floats are an integer part optionally followed by a decimal point and at least one digit.

\subsection{Booleans}

Constant booleans are the keywords \term{True} and \term{False}.

\subsection{Timestamps}

A timestamp literal is encoded as \term{T(}\nterm{string}\term{)} where the string is an ISO 8601 conformant timestamp. The format is \texttt{YYYY-MM-DDThh:mm:ss[.sss]I}, where:
\begin{itemize}
    \item \texttt{YYYY-MM-DD} is the date (year, month, day separated by dashes)
    \item \texttt{hh:mm:ss} is the time (hours, minutes, seconds separated by colons)
    \item \texttt{.sss} is an optional fractional seconds component
    \item \texttt{I} is timezone information: either \texttt{Z} for UTC or an offset $\pm$\texttt{hh:mm}
\end{itemize}

\subsection{Durations}

A duration literal is encoded as \term{D(}\nterm{int}\term{,} \nterm{int}\term{,} \nterm{int}\term{,} \nterm{int} | \nterm{float}\term{)} representing days, hours, minutes, and seconds respectively. Values exceeding their natural bounds overflow into the next larger unit (e.g., 90 seconds becomes 1 minute 30 seconds).

\subsection{None}

The keyword \term{None} represents the absence of a value (the bottom element $\bot$ of partial functions).

\subsection{Symbolic Names as Values}

A bare symbolic name in an expression context refers to a value bound to that name in the current record. The string representation of an expression without an explicit alias (via \term{AS}) is used as the column name through a \term{ValueReference} mechanism.

\section{Records and Tables}

OPQL evaluation is structured around a current state consisting of an event log $L$ and a table $\mathcal{B}$. A table is a multiset of records, each carrying the bindings established by preceding clauses. Tables accumulate and transform as the query executes: PATTERN extends records with entity bindings; KEEP and WHEN compute new bound values; FILTER modifies the underlying event log. Query evaluation begins with the table $[\{ \}]$ containing a single empty record.

\subsection{Records}

A record is a finite mapping from symbolic names to values, representing one row of a table. Each symbolic name in a record corresponds to a column.

\begin{frameddef}[Records]
Let $\mathbb{U}_{sym}$ be the universe of symbolic names and $\mathbb{U}_V$ the universe of values. A \emph{record} is a set of tuples $b \in \mathbb{P}(\mathbb{U}_{sym} \times \mathbb{U}_V)$, where each element associates a concrete value to a symbolic name. We refer to $\mathbb{U}_B = \mathbb{P}(\mathbb{U}_{sym} \times \mathbb{U}_V)$ as the universe of records.

The function $c : \mathbb{U}_B \times \mathbb{U}_{sym} \nrightarrow \mathbb{U}_V$ returns the value from a record $b$ bound to symbolic name $a$:
\[
c(b, a) = v \text{ iff } (a, v) \in b.
\]

The concatenation $\sqcup : \mathbb{U}_B \times \mathbb{U}_B \nrightarrow \mathbb{U}_B$ of two records with distinct symbolic names is:
\[
b \sqcup b' = b \cup b' \text{ s.t. } \not\exists_{(a,v) \in b,\, (a',v') \in b'}\ a = a',
\]
i.e., no duplicate symbolic names are allowed in $b \sqcup b'$.
\end{frameddef}

Record concatenation is partial: it is undefined when both records bind the same symbolic name. This prevents earlier bindings from being silently overwritten; a symbolic name established by a PATTERN clause cannot be rebound by a subsequent KEEP clause using the same name.

\subsection{Tables}

\begin{frameddef}[Tables]
A \emph{multiset} $\mathcal{B} = (\mathbb{U}_B, m)$ over the universe of records is called a \emph{table}. We refer to the set of all possible tables $\mathbb{U}_T = \{ (\mathbb{U}_B, m) \mid m : \mathbb{U}_B \to \mathbb{N} \}$ as the universe of tables.
\end{frameddef}

Tables are multisets rather than sets: the same record may appear multiple times, and each occurrence is processed independently by subsequent clauses. Duplicate records arise naturally when multiple entity bindings satisfy the same pattern; they are preserved unless \term{DISTINCT} is used to collapse them.

\section{Patterns}

Patterns define structural constraints over subgraphs of events, objects, and their qualified relations. A PATTERN clause matches all subgraphs of the event log that satisfy the pattern and binds the matched entities to symbolic names, producing one new record extension per match.

\subsection{Events and Objects}

Events and objects are the two entity types in an OCEL. A pattern node represents one such entity: it may carry a symbolic name (by which the entity is referenced in subsequent clauses), a type qualifier (constraining which event or object types are accepted), or both, or neither.

\begin{frameddef}[Events, Objects]
A tuple $v$ defines how an entity shall be queried:
\[
v = (t, s, q)
\]
with $t \in \{ event, object \}$, $s \in \mathbb{U}_{sym} \cup \{ None \}$, $q \in \Sigma^* \cup \{ None \}$.

The elements are referred to as $\pi_t(v)$, $\pi_s(v)$, and $\pi_q(v)$. The universe of entities is $\mathbb{U}_{entity} = \{ event, object \} \times (\mathbb{U}_{sym} \cup \{ None \}) \times (\Sigma^* \cup \{ None \})$.
\end{frameddef}

\begin{framedgram}[Event, Object]
\begin{tabularx}{\linewidth}{l c X}
\nterm{event} & = & \term{E(} [ \nterm{symbolicName} ] [ \term{:} \nterm{string} ] \term{)} \\
\nterm{object} & = & \term{O(} [ \nterm{symbolicName} ] [ \term{:} \nterm{string} ] \term{)} \\
\end{tabularx}
\end{framedgram}

The optional symbolic name encodes $s$; the optional string literal (preceded by \term{:}) encodes $q$, the required event or object type. The value of $t$ is determined by whether \nterm{event} or \nterm{object} is used. If no symbolic name is given, a globally unique implementation-defined name is assigned so that the entity participates in relation constraints without being accessible by name.

\subsection{Relations}

Relations connect pattern nodes. Three directional forms are available, corresponding to OCEL event-to-object and object-to-object relations, with direction indicating which entity in the pair is considered the source.

\begin{frameddef}[Relations]
A relation is a tuple $r = (d, s, q)$ with $d \in \{ \rightarrow, \leftarrow, \rightleftarrows \}$, $s \in \mathbb{U}_{sym} \cup \{ None \}$, $q \in \Sigma^* \cup \{ None \}$.

The universe of relations is $\mathbb{U}_{relation} = \{ \rightarrow, \leftarrow, \rightleftarrows \} \times (\mathbb{U}_{sym} \cup \{ None \}) \times (\Sigma^* \cup \{ None \})$. Elements are referred to as $\pi_d(r)$, $\pi_s(r)$, and $\pi_q(r)$.
\end{frameddef}

\begin{framedgram}[Relations]
\begin{tabularx}{\linewidth}{l c X}
\nterm{relationLeft} & = & \term{<-[} [ \nterm{symbolicName} ] [ \term{:} \nterm{string} ] \term{]-} \\
\nterm{relationRight} & = & \term{-[} [ \nterm{symbolicName} ] [ \term{:} \nterm{string} ] \term{]->} \\
\nterm{relationAny} & = & \term{-[} [ \nterm{symbolicName} ] [ \term{:} \nterm{string} ] \term{]-} \\
\end{tabularx}
\end{framedgram}

\nterm{relationLeft} encodes $\leftarrow$, \nterm{relationRight} encodes $\rightarrow$, and \nterm{relationAny} encodes $\rightleftarrows$. The direction of event-to-object relations indicates which side carries the event: $\rightarrow$ means the left node is the event; $\leftarrow$ means the right node is the event. Object-to-object relations use $\rightarrow$ and $\leftarrow$ to express directed O2O relationships, and $\rightleftarrows$ to match either direction. The optional string after \term{:} constrains the relation qualifier $q$.

\subsection{Pattern Composition}

A pattern is a chain of alternating entities and relations. The grammar restricts which combinations are valid: object-to-event direction is not permitted in event-to-object connections, and direct event-to-event connections are not permitted at all (events interact only through shared objects).

\begin{frameddef}[Patterns]
A pattern is a finite, alternating sequence
\[
\sigma = \langle v_1, r_1, v_2, r_2, \dots, r_{n-1}, v_n \rangle,\ n \in \mathbb{N},\ v_i \in \mathbb{U}_{entity},\ r_i \in \mathbb{U}_{relation}
\]
such that:
\[
\forall_{i \in \{1,\dots,n\}}\ v_i r_i v_{i+1} \neq (object, s, q)(\rightarrow, s', q')(event, s'', q'')
\]
\[
\land\ \forall_{i \in \{1,\dots,n\}}\ v_i r_i v_{i+1} \neq (event, s, q)(\leftarrow, s', q')(object, s'', q'')),
\]
i.e., event-to-object relations never directionally point from objects to events, and no event-to-event relations are permitted. The universe of all such patterns is $\mathbb{U}_{pattern}$.
\end{frameddef}

\begin{framedgram}[Pattern]
\begin{tabularx}{\linewidth}{l c X}
\nterm{patternObject} & = & \nterm{object} \textbar\ \\
& & (\nterm{object}, (\nterm{relationLeft} \textbar\ \nterm{relationAny}), \nterm{pattern}) \textbar\ \\
& & (\nterm{object}, \nterm{relationRight}, \nterm{patternObject}) \\
\nterm{pattern} & = & \nterm{event} \textbar\ \nterm{object} \textbar\ \\
& & (\nterm{event}, (\nterm{relationRight} \textbar\ \nterm{relationAny}), \nterm{patternObject}) \textbar\ \\
& & (\nterm{object}, (\nterm{relationAny} \textbar\ \nterm{relationLeft}), \nterm{pattern}) \textbar\ \\
& & (\nterm{object}, \nterm{relationRight}, \nterm{patternObject}) \\
\end{tabularx}
\end{framedgram}

The grammar encodes the sequence structure as a recursive right-expansion. Each pattern starts with a \nterm{pattern} node, which may be an event or object optionally followed by a relation and further nodes. Patterns starting with an event may only continue into \nterm{patternObject} chains (no events reachable from an event); patterns starting with an object may continue to either another object or an event.

\subsection{Pattern Constraints}

Given a log $L$ and a record $b$, a pattern imposes three classes of constraint on the values bound to its symbolic names: type constraints on individual entities, directional constraints on event-to-object relations, and directional constraints on object-to-object relations.

\begin{frameddef}[Entity Constraints]
Let $L \in \mathbb{U}_{OCEL}$ and $\sigma \in \mathbb{U}_{pattern}$. The constraints that entities $v \in \mathbb{U}_{entity}$ in $\sigma$ impose on values $c(b, \pi_s(v))$ bound to symbolic names in a record $b \in \mathbb{U}_B$ are given by $R_{EO} : \mathbb{U}_{OCEL} \times \mathbb{U}_B \times \mathbb{U}_{pattern} \to \mathbb{U}_{bool}$:
\[
\begin{array}{rcrl}
R_{EO}(L, b, \sigma) = & & \displaystyle \bigwedge_{i \in Ev(\sigma)} & c(b, \pi_s(v_i)) \in E_L \\
& \land & \displaystyle \bigwedge_{i \in Ev(\sigma)} & (\pi_q(v_i) = None \oplus \pi_q(v_i) = evtype(c(b, \pi_s(v_i)))) \\
& \land & \displaystyle \bigwedge_{i \in Ob(\sigma)} & c(b, \pi_s(v_i)) \in O_L \\
& \land & \displaystyle \bigwedge_{i \in Ob(\sigma)} & (\pi_q(v_i) = None \oplus \pi_q(v_i) = objtype(c(b, \pi_s(v_i))))
\end{array}
\]
with $Ev(\sigma) = \{ i \mid v_i \in \sigma \land \pi_t(v_i) = event \}$ and $Ob(\sigma) = \{ i \mid v_i \in \sigma \land \pi_t(v_i) = object \}$.
\end{frameddef}

Each event node must be bound to an event in $E_L$ and, if a type qualifier is given, the event's type must match. Object nodes are constrained analogously. The XOR $\oplus$ in the type qualifier condition encodes that either no qualifier is present or the qualifier matches (but not both simultaneously being absent/mismatched).

\begin{frameddef}[Event-to-Object Constraints]
Let $L \in \mathbb{U}_{OCEL}$, $b \in \mathbb{U}_B$, $\sigma \in \mathbb{U}_{pattern}$. Event-to-object relations constrain records via $R_{E2O} : \mathbb{U}_{OCEL} \times \mathbb{U}_B \times \mathbb{U}_{pattern} \to \mathbb{U}_{bool}$:
\[
\begin{array}{rrrl}
R_{E2O}(L, b, \sigma) = & & \displaystyle \bigwedge_{i \in E2O_{\leftarrow}(\sigma)} & (c(b, \pi_s(v_i)), c(b, \pi_s(r_i)), c(b, \pi_s(v_{i+1}))) \in E2O_L \\
& \land & \displaystyle \bigwedge_{i \in E2O_{\rightarrow}(\sigma)} & (c(b, \pi_s(v_{i+1})), c(b, \pi_s(r_i)), c(b, \pi_s(v_i))) \in E2O_L \\
& \land & \displaystyle \bigwedge_{i \in E2O(\sigma)} & (\pi_q(r_i) = None \oplus \pi_q(r_i) = c(b, \pi_s(r_i)))
\end{array}
\]
with $E2O(\sigma) = \{ i \mid v_i \in \sigma \land \pi_t(v_i) = event \} \cup \{ i \mid v_{i+1} \in \sigma \land \pi_t(v_{i+1}) = event \}$, $E2O_{\leftarrow}(\sigma) = \{ i \in E2O(\sigma) \mid \pi_t(v_i) = event \}$, $E2O_{\rightarrow}(\sigma) = \{ i \in E2O(\sigma) \mid \pi_t(v_{i+1}) = event \}$.
\end{frameddef}

For a \nterm{relationRight} edge pointing from event to object, the E2O triple must exist in $E2O_L$ with the event as the first element. For a \nterm{relationLeft} edge (object side on the left of the syntax, event on the right), the triple is reversed accordingly. The qualifier constraint uses the same XOR logic as entity type constraints.

\begin{frameddef}[Object-to-Object Constraints]
Let $L \in \mathbb{U}_{OCEL}$, $b \in \mathbb{U}_B$, $\sigma \in \mathbb{U}_{pattern}$. Object-to-object relations constrain records via $R_{O2O} : \mathbb{U}_{OCEL} \times \mathbb{U}_B \times \mathbb{U}_{pattern} \to \mathbb{U}_{bool}$:
\[
\begin{array}{rrrl}
R_{O2O}(L, b, \sigma) = & & \displaystyle \bigwedge_{i \in O2O_{\leftarrow}(\sigma)} & (c(b, \pi_s(v_i)), c(b, \pi_s(r_i)), c(b, \pi_s(v_{i+1}))) \in O2O_L \\
& \land & \displaystyle \bigwedge_{i \in O2O_{\rightarrow}(\sigma)} & (c(b, \pi_s(v_{i+1})), c(b, \pi_s(r_i)), c(b, \pi_s(v_i))) \in O2O_L \\
& \land & \displaystyle \bigwedge_{i \in O2O_{\rightleftarrows}(\sigma)} & \begin{gathered} (c(b, \pi_s(v_i)), c(b, \pi_s(r_i)), c(b, \pi_s(v_{i+1}))) \in O2O_L \\ \lor\ (c(b, \pi_s(v_{i+1})), c(b, \pi_s(r_i)), c(b, \pi_s(v_i))) \in O2O_L \end{gathered} \\
& \land & \displaystyle \bigwedge_{i \in O2O(\sigma)} & (\pi_q(r_i) = None \oplus \pi_q(r_i) = c(b, \pi_s(r_i)))
\end{array}
\]
with $O2O(\sigma) = \{ i \mid v_i, v_{i+1} \in \sigma \land \pi_t(v_i) = \pi_t(v_{i+1}) = object \}$, $O2O_{\leftarrow}(\sigma) = \{ i \in O2O(\sigma) \mid \pi_d(r_i) = \leftarrow \}$, $O2O_{\rightarrow}(\sigma) = \{ i \in O2O(\sigma) \mid \pi_d(r_i) = \rightarrow \}$, $O2O_{\rightleftarrows}(\sigma) = \{ i \in O2O(\sigma) \mid \pi_d(r_i) = \rightleftarrows \}$.
\end{frameddef}

The undirected relation $\rightleftarrows$ is satisfied if the O2O triple exists in either direction. Directed relations $\rightarrow$ and $\leftarrow$ require the triple to exist with the specified orientation.

\begin{frameddef}[Pattern Constraints]
Let $L \in \mathbb{U}_{OCEL}$. A sequence $\sigma \in \mathbb{U}_{pattern}$ encodes the criteria $R(L, b, \sigma)$ any record $b \in \mathbb{U}_B$ must fulfill:
\[
R(L, b, \sigma) = R_{EO}(L, b, \sigma) \land R_{E2O}(L, b, \sigma) \land R_{O2O}(L, b, \sigma).
\]
\end{frameddef}

\subsection{PATTERN Clause}

A PATTERN clause specifies one or more patterns that must all be satisfied simultaneously by the same record extension. Each combination of OCEL entities that satisfies all patterns and is compatible with the existing record produces one new record in the output table.

\begin{framedgram}[PATTERN Clause]
\begin{tabularx}{\linewidth}{l c X}
\nterm{patternList} & = & \nterm{pattern} \{ \term{,} \nterm{pattern} \} \\
\nterm{patternClause} & = & \term{PATTERN} \nterm{patternList} [ \term{SUBJECTTO} | \term{ST} \nterm{expression} ] \\
\end{tabularx}
\end{framedgram}

\begin{frameddef}[PATTERN Clause Evaluation]
The universe of PATTERN clauses is $\mathbb{U}_{pC} = \mathbb{P}(\mathbb{U}_{pattern})$. The function $eval_{pC} : \mathbb{U}_{OCEL} \times \mathbb{U}_T \times \mathbb{U}_{pC} \to \mathbb{U}_T$ evaluates a PATTERN clause:
\[
eval_{pC}(L, \mathcal{B}, \{ \sigma_1, \dots, \sigma_n \}) = \biguplus_{b \in \mathcal{B}} [\, b \sqcup b' \mid R(L, b \sqcup b', \sigma_i)\ \forall i \in \{1 \dots n\},\ b' \in \mathbb{U}_B \,].
\]
\end{frameddef}

For each existing record $b$ in the table, every binding $b'$ that extends $b$ to satisfy all $n$ patterns simultaneously is added. One input record may expand into many output records. Patterns in the list are evaluated conjunctively: a single binding $b'$ must satisfy all of them, not just one.

\section{Expressions}

Expressions compute values from the current record. They may reference bound symbolic names, constant literals, event and object properties, and function calls. The expression language is stratified: arithmetic expressions form the base, comparisons build on them, and propositional connectives combine comparisons.

\subsection{Value References}

A bare symbolic name in an expression refers to the value bound to that name in the current record. A property reference accesses a named attribute of an event or object. Object attributes are versioned: their value depends on the timestamp at which they are evaluated, and a timestamp must be supplied via \term{@}. Event attributes are not versioned and do not require a timestamp.

\begin{framedgram}[Value References]
\begin{tabularx}{\linewidth}{l c X}
\nterm{propertyTimestamp} & = & \nterm{symbolicName} | \nterm{timestamp} \\
\nterm{property} & = & \nterm{symbolicName} \term{[} \nterm{string} [ \term{@} \nterm{propertyTimestamp} ] \term{]} \\
\end{tabularx}
\end{framedgram}

The \term{@} qualifier may be a symbolic name bound to an event (whose event time is then used) or a literal timestamp. Accessing an object attribute without a \term{@} qualifier is undefined, except for the special attributes \texttt{"ocel\_type"} and \texttt{"ocel\_id"}.

\begin{frameddef}[Evaluation of Properties]
Let $L \in \mathbb{U}_{OCEL}$, $b \in \mathbb{U}_B$. A property encodes a tuple $(e, a, t) \in \mathbb{U}_{prop}$ with universe $\mathbb{U}_{prop} = \mathbb{U}_{sym} \times \mathbb{U}_\Sigma \times (\mathbb{U}_{sym} \cup \mathbb{U}_{time} \cup \{ None \})$.

The evaluation $eval_{prop} : \mathbb{U}_{OCEL} \times \mathbb{U}_B \times \mathbb{U}_{prop} \nrightarrow \mathbb{U}_V$ dispatches on whether the symbolic name is bound to an event or an object:
\[
eval_{prop}(L, b, (e, a, t)) = \begin{cases}
eval_{eprop}(L, b, (e, a, t)) & \text{if } c(b, e) \in E_L \\
eval_{oprop}(L, b, (e, a, t)) & \text{if } c(b, e) \in O_L
\end{cases}
\]
with $eval_{eprop} : \mathbb{U}_{OCEL} \times \mathbb{U}_B \times \mathbb{U}_{prop} \nrightarrow \mathbb{U}_V$:
\[
eval_{eprop}(L, b, (e, a, t)) = \begin{cases}
evtype(c(b, e)) & \text{if } a = \text{"ocel\_type"} \\
time(c(b, e)) & \text{if } a = \text{"ocel\_time"} \\
eaval_a(c(b, e)) & \text{otherwise}
\end{cases}
\]
and $eval_{oprop} : \mathbb{U}_{OCEL} \times \mathbb{U}_B \times \mathbb{U}_{prop} \nrightarrow \mathbb{U}_V$:
\[
eval_{oprop}(L, b, (e, a, t)) = \begin{cases}
objtype(c(b, e)) & \text{if } a = \text{"ocel\_type"} \\
oaval_a^t(c(b, e)) & \text{if } t \in \mathbb{U}_{time} \\
oaval_a^{resolve_t(L, b, t)}(c(b, e)) & \text{if } t \in \mathbb{U}_{sym}
\end{cases}
\]
where $resolve_t$ resolves a symbolic timestamp qualifier:
\[
resolve_t(L, b, t) = \begin{cases}
c(b, t) & \text{if } c(b, t) \in \mathbb{U}_{time} \\
time(c(b, t)) & \text{if } c(b, t) \in E_L
\end{cases}
\]
\end{frameddef}

The special attribute names \texttt{"ocel\_type"} and \texttt{"ocel\_time"} are handled by the evaluation functions rather than looked up in the attribute tables. The \term{@} qualifier resolves to a timestamp in one of two ways: if the symbolic name is bound to a timestamp value (e.g., one computed by a prior KEEP clause), that timestamp is used directly; if the symbolic name is bound to an event, the event's \texttt{ocel\_time} is used as shorthand --- this is syntactic sugar equivalent to explicitly writing the event's timestamp.

\subsection{Subquery Expressions}

A subquery expression embeds a full query as a value. The query is enclosed in parentheses and may be prefixed with a materialization hint. Subquery expressions are primarily useful in KEEP clauses, where the result is bound to a symbolic name and can be referenced by subsequent clauses or the RETURN statement.

\begin{framedgram}[Subquery Expressions]
\begin{tabularx}{\linewidth}{l c X}
\nterm{subquery} & = & [ \term{MATERIALIZED} ] \term{(} \nterm{fullQuery} \term{)} \\
                 & | & \term{NOT MATERIALIZED} \term{(} \nterm{fullQuery} \term{)} \\
\end{tabularx}
\end{framedgram}

A subquery expression is unambiguous with the grouping parentheses of arithmetic expressions (\term{(} \nterm{expression} \term{)}), because a \nterm{fullQuery} begins with a clause keyword (\term{PATTERN}, \term{FILTER}, \term{KEEP}, \term{WHEN}) or \term{RETURN}, none of which are valid expression starts.

\begin{frameddef}[Materialization]
A subquery expression carries a materialization mode $m \in \{ eager, lazy \}$:
\begin{itemize}
    \item \textbf{Eager} (default, or explicit \term{MATERIALIZED}): the subquery is evaluated immediately when the enclosing expression is evaluated. The result is a table $\mathcal{B} \in \mathbb{U}_T$.
    \item \textbf{Lazy} (\term{NOT MATERIALIZED}): the subquery is stored as an unevaluated AST node. It is evaluated each time the bound symbolic name is referenced.
\end{itemize}

Let $Q \in \mathbb{U}_{opql}$ be the embedded query. The evaluation $eval_{sq} : \mathbb{U}_{OCEL} \times \mathbb{U}_B \times \mathbb{U}_{opql} \times \{eager, lazy\} \nrightarrow \mathbb{U}_T \cup \mathbb{U}_{opql}$ is:
\[
eval_{sq}(L, b, Q, m) = \begin{cases}
eval_{opql}(L, [\{\}], Q) & \text{if } m = eager \\
Q & \text{if } m = lazy
\end{cases}
\]

When a lazy subquery is later referenced via a symbolic name lookup and the bound value is an unevaluated query $Q$, evaluation proceeds as $eval_{opql}(L, [\{\}], Q)$.
\end{frameddef}

An eager subquery is evaluated once and its result stored; referencing the bound name multiple times does not re-execute the query. A lazy subquery is re-executed on each reference.

\paragraph{Note on current equivalence.} Under the current language, the materialization mode has no observable effect on query results. The only clause that mutates the event log is FILTER, but FILTER also clears all context bindings (see Section~9), so a lazy subquery can never be referenced after the log it would execute against has changed. Since OPQL does not permit rebinding of symbolic names, the inputs to a lazy subquery are always identical across references within the same scope. The distinction is retained in anticipation of future log-editing operations (e.g., clauses that add events, objects, or relations to the log without clearing context), at which point a lazy subquery would re-execute against the modified log while an eager subquery would retain its original result. Whether FILTER's full context reset remains the appropriate design once such operations exist is an open question.

\paragraph{Undefined usage.} Subquery expressions evaluate to tables, not scalar values. This standard defines their behaviour only when consumed by aggregation functions or bound via KEEP for later aggregation. Using a table-valued subquery expression in a context that expects a scalar (e.g., as a bare RETURN column, as an operand in arithmetic, or in a comparison) is undefined. The behaviour of the reference implementation in such cases is not authoritative.

\subsection{Arithmetic Expressions}

Arithmetic expressions compute numeric, temporal, and duration values. Unary \term{+} and \term{-} apply to a single operand; all binary operators apply to two operands.

\begin{framedgram}[Arithmetic Expressions]
\begin{tabularx}{\linewidth}{l c X}
\nterm{value} & = & \nterm{symbolicName} | \nterm{constValue} | \nterm{property} | \nterm{functionCall} | \nterm{subquery} \\
\nterm{arExpr} & = & \nterm{value} | \\
& & \term{+} \nterm{arExpr} | \term{-} \nterm{arExpr} | \\
& & \nterm{arExpr} \term{+} \nterm{arExpr} | \nterm{arExpr} \term{-} \nterm{arExpr} | \\
& & \nterm{arExpr} \term{*} \nterm{arExpr} | \nterm{arExpr} \term{/} \nterm{arExpr} | \nterm{arExpr} \term{\%} \nterm{arExpr} | \\
& & \nterm{arExpr} \term{\textasciicircum} \nterm{arExpr} \\
\end{tabularx}
\end{framedgram}

\begin{frameddef}[Arithmetic Operations]
The operations $\mathbb{U}_{ops} = \{ +, -, *, /, \%, \wedge \}$ are defined on:
\begin{center}
\begin{tabular}{ll}
$+ : \mathbb{R} \times \mathbb{R} \to \mathbb{R}$ & $- : \mathbb{R} \times \mathbb{R} \to \mathbb{R}$ \\
$+ : \mathbb{U}_{time} \times \mathbb{U}_{duration} \to \mathbb{U}_{time}$ & $- : \mathbb{U}_{time} \times \mathbb{U}_{time} \to \mathbb{U}_{duration}$ \\
$+ : \mathbb{U}_{duration} \times \mathbb{U}_{time} \to \mathbb{U}_{time}$ & $- : \mathbb{U}_{time} \times \mathbb{U}_{duration} \to \mathbb{U}_{time}$ \\
$* : \mathbb{R} \times \mathbb{R} \to \mathbb{R}$ & $/ : \mathbb{R} \times \mathbb{R} \to \mathbb{R}$ \\
$* : \mathbb{U}_{duration} \times \mathbb{R} \to \mathbb{U}_{duration}$ & $/ : \mathbb{U}_{duration} \times \mathbb{R} \to \mathbb{U}_{duration}$ \\
$* : \mathbb{R} \times \mathbb{U}_{duration} \to \mathbb{U}_{duration}$ & $/ : \mathbb{U}_{duration} \times \mathbb{U}_{duration} \to \mathbb{R}$ \\
$\% : \mathbb{N} \times \mathbb{N} \to \mathbb{N}$ & $\wedge : \mathbb{R} \times \mathbb{R} \to \mathbb{R}$
\end{tabular}
\end{center}
Binding priority $bp : \mathbb{U}_{ops} \to \mathbb{N}$: $\{ (\wedge, 3), (*, 2), (/, 2), (\%, 2), (+, 1), (-, 1) \}$. Operators with equal binding priority associate left-to-right; $\wedge$ is right-associative.
\end{frameddef}

An operation applied to operands not listed in the type table above is undefined and yields $None$. Any arithmetic operation with a $None$ operand yields $None$. Timestamps may be shifted by adding or subtracting a duration; subtracting two timestamps yields the duration between them. Durations may be scaled by a real number, and dividing two durations yields a real number.

\subsection{Comparisons}

Comparisons produce boolean values by comparing two arithmetic expressions. They are building blocks for propositional logic expressions.

\begin{frameddef}[Comparisons]
The universe of relations is $\mathbb{U}_{rel} = \{ =, \leq, \geq, <, >, \neq \}$. Relations are functions $\circ : V \times V \to \mathbb{U}_{bool}$ for $V \in \{ \mathbb{R}, \mathbb{U}_{time}, \mathbb{U}_{duration}, \mathbb{U}_{bool}, \mathbb{U}_\Sigma \}$.

A comparison expression $(c_1, \circ, c_2)$ with $c_1, c_2 \in \mathbb{U}_{arEx}$, $\circ \in \mathbb{U}_{rel} \cup \{ None \}$ evaluates as:
\[
eval_{comp}(L, b, (c_1, \circ, c_2)) = \begin{cases}
eval_{arEx}(L, b, c_1) \circ eval_{arEx}(L, b, c_2) & \text{if } \circ \neq None \\
eval_{arEx}(L, b, c_1) & \text{if } \circ = c_2 = None
\end{cases}
\]
\end{frameddef}

The second case ($\circ = None$) occurs when a comparison expression reduces to a bare arithmetic expression — for example, a bare boolean value or a function returning a boolean. Comparing values of incompatible types is undefined and yields $None$. Any comparison involving a $None$ operand yields $None$.

\subsection{Propositional Logic}

Propositional logic expressions combine comparisons using the three boolean connectives \term{AND}, \term{XOR}, and \term{OR}, and the unary \term{NOT}. A \emph{logical atom} $l \in \mathbb{U}_{logA}$ is either a comparison expression or a comparison expression preceded by \term{NOT}.

\begin{framedgram}[Logical Expressions]
\begin{tabularx}{\linewidth}{l c X}
\nterm{expr} & = & \nterm{cExpr} | \term{NOT} \nterm{cExpr} | \\
& & \nterm{expr} \term{AND} \nterm{expr} | \nterm{expr} \term{XOR} \nterm{expr} | \nterm{expr} \term{OR} \nterm{expr} \\
\end{tabularx}
\end{framedgram}

\begin{frameddef}[Logical Connectives]
The universe of logical connectives is $\mathbb{U}_{logCon} = \{ \land, \lor, \oplus \} \subset (\mathbb{U}_{bool} \times \mathbb{U}_{bool} \to \mathbb{U}_{bool})$. Binding priority $bp : \mathbb{U}_{logCon} \to \mathbb{N}$: $\{ (\land, 3), (\oplus, 2), (\lor, 1) \}$.
\end{frameddef}

\term{AND} binds more tightly than \term{XOR}, which binds more tightly than \term{OR}. All three connectives are left-associative (chains of the same connective associate left-to-right).

\begin{frameddef}[Logical Atom Evaluation]
The evaluation $eval_{logA} : \mathbb{U}_{OCEL} \times \mathbb{U}_B \times \mathbb{U}_{logA} \nrightarrow \mathbb{U}_{bool}$ is:
\[
eval_{logA}(L, b, l) = \begin{cases}
eval_{comp}(L, b, c) & \text{if } l = c \in \mathbb{U}_{comp} \\
\lnot\, eval_{comp}(L, b, c) & \text{if } l = (\text{NOT},\ c),\ c \in \mathbb{U}_{comp}
\end{cases}
\]
\end{frameddef}

\begin{frameddef}[Logical Expression Evaluation]
Grammar encodes logical expressions as sequences $\tau = \langle l_1, \circ_1, l_2, \dots, \circ_{n-1}, l_n \rangle$ with $l_i \in \mathbb{U}_{logA}$, $\circ_i \in \mathbb{U}_{logCon}$.

The evaluation $eval_{lExpr} : \mathbb{U}_{OCEL} \times \mathbb{U}_B \times \mathbb{U}_{lExpr} \nrightarrow \mathbb{U}_V$ proceeds recursively, pivoting on the rightmost operator with the lowest binding priority:
\[
eval_{lExpr}(L, b, \tau) = \begin{cases}
eval_{logA}(L, b, l) & \text{if } \tau = \langle l \rangle \\
eval_{lExpr}(L, b, \tau_L) \circ_i\, eval_{lExpr}(L, b, \tau_R) & \text{otherwise}
\end{cases}
\]
where $i = \max(\{ j \mid \not\exists k \neq j : bp(\circ_k) < bp(\circ_j) \})$, $\tau_L = \langle l_1, \dots, \circ_{i-1}, l_i \rangle$, and $\tau_R = \langle l_{i+1}, \dots, \circ_{n-1}, l_n \rangle$.
\end{frameddef}

The pivot $i$ selects the rightmost operator that is outermost in the expression tree (lowest binding priority = loosest binding = last evaluated). This produces standard left-to-right evaluation for chains of equal-priority operators: $A \text{ AND } B \text{ AND } C$ evaluates as $(A \text{ AND } B) \text{ AND } C$. The \term{AND}-before-\term{XOR}-before-\term{OR} precedence mirrors standard boolean algebra.

\section{Binning}

Binning maps a numeric value to a user-defined target value based on which of a set of intervals it falls into. It is applied as a modifier to a projection item, replacing the computed value with the label of the matching bin.

\begin{framedgram}[Binning]
\begin{tabularx}{\linewidth}{l c X}
\nterm{intervalTarget} & = & \nterm{int} | \nterm{float} | \nterm{string} \\
\nterm{intervalLimit} & = & \nterm{int} | \nterm{float} | \term{-inf} | [ \term{+} ] \term{inf} \\
\nterm{binningInterval} & = & (\term{(} | \term{[}) \nterm{intervalLimit} \term{,} \nterm{intervalLimit} \\
& & \quad (\term{)} | \term{]}) \term{AS} \nterm{intervalTarget} \\
\nterm{binning} & = & \term{BINNED} \term{(} \nterm{binningInterval} \{ \term{,} \nterm{binningInterval} \} \term{)} \\
\end{tabularx}
\end{framedgram}

Each interval is written as a standard mathematical interval. A round bracket \term{(} or \term{)} denotes an open (exclusive) bound; a square bracket \term{[} or \term{]} denotes a closed (inclusive) bound. The limits \term{-inf} and \term{inf} represent $-\infty$ and $+\infty$ respectively. The \term{AS} target specifies the value the expression is replaced with when it falls within the interval; the target may be a numeric literal or a string.

\begin{frameddef}[Binning Interval]
Let $\mathbb{U}_{\infty} = \{ \infty, -\infty \}$ with $\forall_{x \in \mathbb{R}}\ -\infty < x < \infty$. Let $\mathbb{U}_{bounds} = (\mathbb{R} \cup \mathbb{U}_{\infty}) \times \{ <, \leq \}$ and $\mathbb{U}_{intv} = \mathbb{U}_{bounds} \times \mathbb{U}_{bounds}$.

The function $ininterval : \mathbb{R} \times \mathbb{U}_{bounds} \times \mathbb{U}_{bounds} \to \mathbb{U}_{bool}$ determines whether a value is within bounds:
\[
ininterval(v, ((n, \circ), (m, \circ'))) = \begin{cases}
True & \text{if } n \circ v \circ' m \\
False & \text{otherwise}
\end{cases}
\]
\end{frameddef}

An open bracket maps to the strict relation $<$; a closed bracket maps to $\leq$. For example, \texttt{[0, 10)} produces bounds $((0, \leq), (10, <))$, matching values $v$ with $0 \leq v < 10$.

\begin{frameddef}[Binning]
The universe of binnings is $\mathbb{U}_{binnings} = (\mathbb{U}_{intv} \times (\mathbb{R} \cup \Sigma^*))^*$. The function $bin : \mathbb{R} \times \mathbb{U}_{binnings} \nrightarrow \mathbb{R} \cup \Sigma^*$ maps a value to the target of the first matching interval:
\[
bin(v, ((i_1, t_1), \dots, (i_n, t_n))) = \begin{cases}
t_j & \text{if } ininterval(v, i_j) \land \not\exists_{k < j}\ ininterval(v, i_k) \\
None & \text{otherwise}
\end{cases}
\]
\end{frameddef}

Intervals are evaluated in the order they appear; the first match determines the result. A value matched by no interval evaluates to $None$. Overlapping intervals are permitted; the first one listed takes precedence. Binning may only be applied to numeric expressions; applying it to a non-numeric value is undefined.

\section{Query Clauses}

\subsection{KEEP Clause}

\begin{framedgram}[KEEP Clause]
\begin{tabularx}{\linewidth}{l c X}
\nterm{sname} & = & \nterm{expression} [ \term{AS} \nterm{symbolicName} ] \\
\nterm{projectionItem} & = & \nterm{sname} [ \nterm{binning} ] \\
\nterm{projection} & = & [ \term{DISTINCT} ] \term{*} [ \nterm{order} ] [ \nterm{limit} ] \\
& | & [ \term{DISTINCT} ] [ \term{*} \term{,} ] \nterm{projectionItem} \{ \term{,} \nterm{projectionItem} \} \\
& & \quad [ \nterm{order} ] [ \nterm{limit} ] \\
\nterm{keepClause} & = & \term{KEEP} \nterm{projection} [ \term{SUBJECTTO} | \term{ST} \nterm{expression} ] \\
\end{tabularx}
\end{framedgram}

A KEEP clause evaluates expressions and binds their results to symbolic names in the current table. The optional \term{AS} clause assigns an explicit alias; when omitted, the column name is derived from the string representation of the expression via ValueReference.

The \term{*} wildcard indicates that previously bound values in each record are retained. Without it, only the newly computed values appear in the result. The \term{DISTINCT} keyword deduplicates the resulting table.

\begin{frameddef}[Deduplication]
Let $\mathcal{B} \in \mathbb{U}_T$ and $S(\mathcal{B}) = \{ b \mid b \in \mathcal{B} \}$. The deduplication function $\epsilon : \mathbb{U}_T \to \mathbb{U}_T$ is:
\[
\epsilon(\mathcal{B}) = [\, b \mid b \in S(\mathcal{B}) \,].
\]
\end{frameddef}

\begin{frameddef}[KEEP Clause Evaluation]
Let $\mathbb{U}_{proI} = \mathbb{U}_{sym} \times \mathbb{U}_{lExpr} \times (\mathbb{U}_{binnings} \cup \{ None \})$ be the universe of projection items. The universe of KEEP clauses is $\mathbb{U}_{kC} = \mathbb{U}_{bool} \times \mathbb{U}_{bool} \times \mathbb{P}(\mathbb{U}_{proI})$, where the first boolean indicates deduplication and the second indicates retention of previous values.

The function $eval_{namedEx} : \mathbb{U}_{OCEL} \times \mathbb{U}_B \times \mathbb{P}(\mathbb{U}_{proI}) \nrightarrow \mathbb{U}_B$ builds a new record:
\begin{align*}
eval_{namedEx}(L, b, P) = & \ \{ (a, eval_{lExpr}(L, b, e)) \mid (a, e, b') \in P \land b' = None \} \\
& \cup \{ (a, bin(eval_{lExpr}(L, b, e), b')) \mid (a, e, b') \in P \land b' \in \mathbb{U}_{binnings} \}
\end{align*}

The evaluation $eval_{kC} : \mathbb{U}_{OCEL} \times \mathbb{U}_T \times \mathbb{U}_{kC} \nrightarrow \mathbb{U}_T$ is:
\[
eval_{kC}(L, \mathcal{B}, (d, h, P)) = \begin{cases}
\biguplus_{b \in \mathcal{B}} [\, b \sqcup eval_{namedEx}(L, b, P) \,] & \text{if } \neg d \land h \\
\epsilon(\biguplus_{b \in \mathcal{B}} [\, b \sqcup eval_{namedEx}(L, b, P) \,]) & \text{if } d \land h \\
\biguplus_{b \in \mathcal{B}} [\, eval_{namedEx}(L, b, P) \,] & \text{if } \neg d \land \neg h \\
\epsilon(\biguplus_{b \in \mathcal{B}} [\, eval_{namedEx}(L, b, P) \,]) & \text{if } d \land \neg h
\end{cases}
\]

All expressions in $P$ are evaluated under the input record $b$, not under records augmented with results of sibling expressions. This allows parallel evaluation.
\end{frameddef}

\subsection{WHEN Clause}

\begin{framedgram}[WHEN Clause]
\begin{tabularx}{\linewidth}{l c X}
\nterm{whenClause} & = & \term{WHEN} \nterm{expression} \term{AS} \nterm{symbolicName} \\
& & \quad [ \term{SUBJECTTO} | \term{ST} \nterm{expression} ] \\
\end{tabularx}
\end{framedgram}

A WHEN clause finds timestamps at which a given expression becomes $True$, and binds each such timestamp to a symbolic name. The search is driven by the version history of object attributes referenced in the expression: the candidate timestamps are exactly the timestamps at which some object attribute mentioned in the expression changes value. For each candidate, the expression is re-evaluated with the symbolic name bound to that timestamp; matching timestamps each produce one output record.

\begin{frameddef}[Candidate Timestamps]
Let $L \in \mathbb{U}_{OCEL}$, $b \in \mathbb{U}_B$, $\tau \in \mathbb{U}_{lExpr}$, $a' \in \mathbb{U}_{sym}$. The function $candidate : \mathbb{U}_{OCEL} \times \mathbb{U}_B \times \mathbb{U}_{lExpr} \times \mathbb{U}_{sym} \nrightarrow \mathbb{P}(\mathbb{U}_{time})$ yields the set of candidate timestamps:
\[
candidate(L, b, \tau, a') = \biguplus_{\rho \in extract_{arEx}(\tau)} \left[ t^* \middle\vert \begin{array}{l} (e, a, t) \in extract_{oa}(L, b, \rho, a') \\ \land (c(b, e), a, t^*) \in \dom(oaval) \end{array} \right]
\]
where $extract_{arEx}(\tau)$ enumerates the arithmetic sub-expressions of $\tau$, and $extract_{oa}(L, b, \rho, a')$ yields the object-attribute-timestamp triples $(e, a, t)$ referenced in $\rho$ that use $a'$ as their timestamp variable.
\end{frameddef}

The candidate set covers every timestamp at which any object attribute accessed via the symbolic name $a'$ has a recorded value. Evaluating the expression at each such timestamp identifies those timestamps at which the expression is $True$.

\begin{frameddef}[WHEN Clause Evaluation]
The universe of WHEN clauses is $\mathbb{U}_{wC} = \mathbb{U}_{lExpr} \times \mathbb{U}_{sym}$. The function $eval_{aT} : \mathbb{U}_{OCEL} \times \mathbb{U}_B \times \mathbb{U}_{wC} \nrightarrow \mathbb{U}_T$ yields timestamps for which the expression evaluates to $True$:
\[
eval_{aT}(L, b, (\tau, a)) = \left[ \{(a, t)\} \middle\vert \begin{array}{l} t \in candidate(L, b, \tau, a) \\ \land eval_{lExpr}(L, b \sqcup \{(a, t)\}) = True \end{array} \right]
\]

The evaluation $eval_{wC} : \mathbb{U}_{OCEL} \times \mathbb{U}_T \times \mathbb{U}_{wC} \nrightarrow \mathbb{U}_T$ is:
\[
eval_{wC}(L, \mathcal{B}, (\tau, a)) = \begin{cases}
\biguplus_{b \in \mathcal{B}} [\, b \sqcup \{(a, None)\} \,] & \text{if } |eval_{aT}(L, b, (\tau, a))| = 0 \\
\biguplus_{b \in \mathcal{B}} [\, b \sqcup \{(a, t)\} \mid t \in eval_{aT}(L, b, (\tau, a)) \,] & \text{otherwise}
\end{cases}
\]

If no timestamp satisfies the expression, the symbolic name is bound to $None$ rather than dropping the record.
\end{frameddef}

One input record may expand into multiple output records — one per matching timestamp. This mirrors the fan-out behavior of PATTERN clauses.

\subsection{FILTER Clause}

\begin{framedgram}[FILTER Clause]
\begin{tabularx}{\linewidth}{l c X}
\nterm{filterClause} & = & \term{FILTER} \nterm{symbolicName} \{ \term{,} \nterm{symbolicName} \} \\
\end{tabularx}
\end{framedgram}

A FILTER clause removes entities from the underlying event log. The symbolic names in the argument list are looked up in each current record; the entities they are bound to are removed from the log, along with all E2O and O2O relations they participate in. After a FILTER clause executes, the table is reset to a single empty record so that subsequent clauses operate on the filtered log from a clean state.

\begin{frameddef}[FILTER Clause Evaluation]
The universe of FILTER clauses is $\mathbb{U}_{fC} = \mathbb{P}(\mathbb{U}_{sym})$. The evaluation $eval_{fC} : \mathbb{U}_{OCEL} \times \mathbb{U}_T \times \mathbb{U}_{fC} \nrightarrow \mathbb{U}_{OCEL}$ yields a filtered event log $L'$:
\[
eval_{fC}(L, \mathcal{B}, A) = L'
\]
with
\begin{multline*}
L' = (E_L \setminus F_E,\ O_L \setminus F_O,\ EA_L,\ OA_L,\ evtype_L,\ time_L,\ objtype_L,\\
\quad eatype_L,\ oatype_L,\ eaval_L,\ oaval_L,\ E2O_L \setminus F_{E2O},\ O2O_L \setminus F_{O2O})
\end{multline*}
where:
\begin{itemize}
    \item $F = \{ c(b, a) \mid \exists_{a \in A, v \in \mathbb{U}_V}\ (a, v) \in b \}$
    \item $F_E = E_L \cap F$, $F_O = O_L \cap F$
    \item $F_{E2O} = \{ (e, q, o) \in E2O_L \mid e \in F_E \lor o \in F_O \}$
    \item $F_{O2O} = \{ (o_1, q, o_2) \in O2O_L \mid o_1 \in F_O \lor o_2 \in F_O \}$
\end{itemize}
\end{frameddef}

\subsection{SUBJECTTO Subclause}

\begin{framedgram}[SUBJECTTO Subclause]
\begin{tabularx}{\linewidth}{l c X}
\nterm{subjectTo} & = & (\term{SUBJECTTO} | \term{ST}) \nterm{expression} \\
\end{tabularx}
\end{framedgram}

A SUBJECTTO subclause is a guard attached to the clause it accompanies. It is evaluated after the clause produces its output table, and retains only records for which the expression is $True$. Records evaluating to $False$, $None$, or any non-boolean value are dropped. SUBJECTTO may be attached to PATTERN, KEEP, WHEN, and RETURN clauses.

\begin{frameddef}[SUBJECTTO Subclause Evaluation]
The universe of SUBJECTTO subclauses is $\mathbb{U}_{stS} = \mathbb{U}_{lExpr}$. The evaluation $eval_{stS} : \mathbb{U}_{OCEL} \times \mathbb{U}_T \times \mathbb{U}_{stS} \nrightarrow \mathbb{U}_T$ is:
\[
eval_{stS}(L, \mathcal{B}, \tau) = [\, b \in \mathcal{B} \mid eval_{lExpr}(L, b, \tau) = True \,].
\]
\end{frameddef}

\section{RETURN Statement}

\begin{framedgram}[RETURN Statement]
\begin{tabularx}{\linewidth}{l c X}
\nterm{returnRule} & = & \term{RETURN} \term{OCEL} \\
& | & \term{RETURN} \nterm{projection} [ \term{SUBJECTTO} | \term{ST} \nterm{expression} ] \\
\end{tabularx}
\end{framedgram}

RETURN statements terminate a query. They come in two forms:

\begin{itemize}
    \item \textbf{RETURN OCEL} --- returns the (potentially filtered) event log as-is
    \item \textbf{Tabular RETURN} --- returns a table of records, using the same projection mechanism as KEEP clauses
\end{itemize}

The tabular form supports the same modifiers as KEEP: \term{DISTINCT} for deduplication, \term{*} for retaining previously bound values, optional \term{AS} aliases (omitting \term{AS} yields a column name from the expression's string representation), binning, ordering, and limiting.

\begin{frameddef}[RETURN Statement Evaluation]
The universe of RETURN statements is $\mathbb{U}_{rS} = \mathbb{U}_{retTab} \cup \mathbb{U}_{retLog}$, where $\mathbb{U}_{retTab} = \mathbb{U}_{kC}$ and $\mathbb{U}_{retLog}$ is the singleton universe of \term{RETURN OCEL} statements.

The evaluation $eval_{rS} : \mathbb{U}_{OCEL} \times \mathbb{U}_T \times \mathbb{U}_{rS} \nrightarrow \mathbb{U}_T \cup \mathbb{U}_{OCEL}$ is:
\[
eval_{rS}(L, \mathcal{B}, R) = \begin{cases}
L & \text{if } R \in \mathbb{U}_{retLog} \\
eval_{kC}(L, \mathcal{B}, R) & \text{if } R \in \mathbb{U}_{retTab}
\end{cases}
\]
\end{frameddef}

\section{Query Composition}

\begin{framedgram}[Complete Query]
\begin{tabularx}{\linewidth}{l c X}
\nterm{contextRule} & = & \nterm{patternClause} | \nterm{filterClause} | \nterm{keepClause} | \nterm{whenClause} \\
\nterm{fullQuery} & = & \{ \nterm{contextRule} \} \nterm{returnRule} \\
\end{tabularx}
\end{framedgram}

A query is a sequence of zero or more context clauses followed by exactly one RETURN statement. Clauses are evaluated sequentially; each clause receives the current state (event log $L$ and table $\mathcal{B}$) and produces a new state for the next clause. PATTERN, KEEP, and WHEN clauses transform only the table; FILTER clauses transform only the event log and reset the table. The RETURN statement terminates evaluation and produces the final result.

\begin{frameddef}[Clause Evaluation]
Let $\mathbb{U}_{bC} = \mathbb{U}_{pC} \cup \mathbb{U}_{kC} \cup \mathbb{U}_{wC}$. The function $eval_{bC} : \mathbb{U}_{OCEL} \times \mathbb{U}_B \times \mathbb{U}_{bC} \nrightarrow \mathbb{U}_T$ evaluates clauses:
\[
eval_{bC}(L, \mathcal{B}, c) = \begin{cases}
eval_{pC}(L, \mathcal{B}, c) & \text{if } c \in \mathbb{U}_{pC} \\
eval_{kC}(L, \mathcal{B}, c) & \text{if } c \in \mathbb{U}_{kC} \\
eval_{wC}(L, \mathcal{B}, c) & \text{if } c \in \mathbb{U}_{wC}
\end{cases}
\]

Let $\mathbb{U}_{cC} = \mathbb{U}_{bC} \times (\mathbb{U}_{stS} \cup \{ None \})$ be the universe of clauses with optional SUBJECTTO subclauses. The evaluation $eval_{cC} : \mathbb{U}_{OCEL} \times \mathbb{U}_B \times \mathbb{U}_{cC} \nrightarrow \mathbb{U}_T$ is:
\[
eval_{cC}(L, \mathcal{B}, (c, s)) = \begin{cases}
eval_{bC}(L, \mathcal{B}, c) & \text{if } s = None \\
eval_{stS}(L, eval_{bC}(L, \mathcal{B}, c), s) & \text{otherwise}
\end{cases}
\]
\end{frameddef}

\begin{frameddef}[Query Evaluation]
The universe of all OPQL queries is $\mathbb{U}_{opql} = (\mathbb{U}_{cC} \cup \mathbb{U}_{fC})^* \times \mathbb{U}_{rS}$.

The function $eval_{opql}$ evaluates a full query $(\langle C_1, \dots, C_n \rangle, R)$:
\[
eval_{opql}(L, \mathcal{B}, (\langle C_1, \dots, C_n \rangle, R)) = \begin{cases}
eval_{rS}(L, \mathcal{B}, R) & \text{if } n = 0 \\
eval_{opql}(L', [\{ \}], (\langle C_2, \dots, C_n \rangle, R)) & \text{if } C_1 \in \mathbb{U}_{fC} \\
eval_{opql}(L, \mathcal{B}', (\langle C_2, \dots, C_n \rangle, R)) & \text{otherwise}
\end{cases}
\]

where $L' = eval_{fC}(L, \mathcal{B}, C_1)$ and $\mathcal{B}' = eval_{cC}(L, \mathcal{B}, C_1)$.

The entry point is $query : \mathbb{U}_{OCEL} \times \mathbb{U}_{opql} \to \mathbb{U}_T \cup \mathbb{U}_{OCEL}$:
\[
query(L, Q) = eval_{opql}(L, [\{ \}], Q).
\]
\end{frameddef}

FILTER clauses reset the context: the filtered event log $L'$ is passed forward and the table is reset to a single empty record $[\{ \}]$. This prevents dangling references to filtered-out entities. All other clauses manipulate only the table $\mathcal{B}$, passing the event log unchanged.

\subsection{Common Table Expressions}

A common table expression (CTE) names the result of a subquery so it can be referenced by symbolic name in subsequent clauses. CTEs are expressed using a KEEP clause with a subquery expression:

\begin{verbatim}
KEEP (PATTERN E(e) RETURN e["cost"]) AS costs
RETURN avg(costs), count(costs), max(costs)
\end{verbatim}

The subquery \texttt{(PATTERN E(e) RETURN e["cost"])} is a subquery expression (see Section~7) that evaluates to a table. The KEEP clause binds this table to the symbolic name \texttt{costs}. Subsequent references to \texttt{costs} --- whether in aggregation functions, other KEEP clauses, or the RETURN statement --- resolve to the stored table.

No special clause syntax is required: a CTE is simply a KEEP projection item whose expression is a subquery. The existing KEEP evaluation rules apply without modification.

\paragraph{Materialization.} By default, or with an explicit \term{MATERIALIZED} prefix, the subquery is evaluated eagerly during KEEP processing. The \term{NOT MATERIALIZED} prefix defers evaluation: the query AST is stored as-is, and evaluated each time the symbolic name is referenced. See Definition~7.2 for formal semantics. Note that under the current language the two modes are observationally equivalent, since no operation can mutate the event log without also clearing context bindings (see the note in Section~7).

\section{Functions}

Functions are named subroutines callable within expressions. A function call takes one or more arguments, each of which is either an expression or a full subquery. When a subquery is passed, it is evaluated starting from the current record as context, allowing it to reference values bound in the enclosing scope.

\begin{framedgram}[Function Calls]
\begin{tabularx}{\linewidth}{l c X}
\nterm{rvFunctionArg} & = & \nterm{expression} | \nterm{fullQuery} \\
\nterm{rvFunctionCall} & = & \nterm{symbolicName} \term{(} \nterm{rvFunctionArg} \{ \term{,} \nterm{rvFunctionArg} \} \term{)} \\
\end{tabularx}
\end{framedgram}

\subsection{Aggregation Functions}

Aggregation functions operate on a table produced by a subquery argument and reduce it to a scalar. The subquery must return a single-column table; the aggregation is performed over the values in that column. Subqueries passed to aggregation functions are evaluated in the context of the current record, so they may reference symbolic names bound by enclosing clauses.

\begin{frameddef}[count]
$count : \mathbb{U}_T \nrightarrow \mathbb{N}$ returns the number of records in the table:
\[
count(\mathcal{B}) = |\mathcal{B}|.
\]
\end{frameddef}

\begin{frameddef}[sum]
$sum : \mathbb{U}_T \nrightarrow \mathbb{R} \cup \mathbb{U}_{duration}$ adds all values in a single-column table:
\[
sum(\mathcal{B}) = \sum_{v \in [\, v \mid (a, v) \in [\, b \in \mathcal{B} \mid |b| = 1 \,] \,]} v.
\]
\end{frameddef}

\begin{frameddef}[min, max]
$min, max : \mathbb{U}_T \nrightarrow \mathbb{R}$ return the minimum and maximum value in a single-column table:
\[
max(\mathcal{B}) = \max(\{ v \mid (a, v) \in \{ b \in \mathcal{B} \mid |b| = 1 \} \})
\]
\[
min(\mathcal{B}) = \min(\{ v \mid (a, v) \in \{ b \in \mathcal{B} \mid |b| = 1 \} \}).
\]
\end{frameddef}

\begin{frameddef}[avg]
$avg : \mathbb{U}_T \nrightarrow \mathbb{R} \cup \mathbb{U}_{duration}$ returns the arithmetic mean of a single-column table:
\[
avg(\mathcal{B}) = \frac{sum(\mathcal{B})}{count(\mathcal{B})}.
\]
\end{frameddef}

\begin{frameddef}[stddev]
$stddev : \mathbb{U}_T \nrightarrow \mathbb{R}$ returns the sample standard deviation of a single-column table:
\[
stddev(\mathcal{B}) = \left( \frac{\sum_{v} (v - avg(\mathcal{B}))^2}{count(\mathcal{B}) - 1} \right)^{0.5}
\]
where $v$ ranges over all values in the single-column table $\mathcal{B}$.
\end{frameddef}

\subsection{Temporal Functions}

Temporal functions navigate the event history of objects. They look up events that are temporally adjacent to a given event within the scope of a shared object.

\begin{frameddef}[olag, olead]
Let $L \in \mathbb{U}_{OCEL}$, $e \in E_L$, $o \in O_L$. The function $laglead(L, e, o, e', \circ)$ with $\circ \in \{ >, < \}$ yields $True$ if event $e'$ is associated with object $o$ (i.e., $(e', q, o) \in E2O_L$ for some qualifier $q$) and $time(e') \circ time(e)$.

$olag(L, b, e, o)$ returns the $ocel\_id$ of the most recent event strictly preceding $e$ that is associated with object $o$, or $None$ if no such event exists. $olead(L, b, e, o)$ returns the $ocel\_id$ of the earliest event strictly succeeding $e$ that is associated with object $o$, or $None$ if no such event exists.

Both functions accept an optional fifth argument $s \in \mathbb{U}_\Sigma$ filtering by event type: $olag(L, b, e, o, s)$ considers only preceding events of type $s$, and $olead(L, b, e, o, s)$ considers only succeeding events of type $s$.

The arguments $e$ and $o$ are symbolic names bound to an event and an object respectively in the current record.
\end{frameddef}

\subsection{Introspection}

\begin{frameddef}[isnone]
$isnone : \mathbb{U}_{OCEL} \times \mathbb{U}_B \times \mathbb{U}_V \to \mathbb{U}_{bool}$ tests whether a value is $None$:
\[
isnone(L, b, v) = \begin{cases}
True & \text{if } v = None \\
False & \text{otherwise}
\end{cases}
\]
\end{frameddef}

$isnone$ is the standard way to test whether a property access, temporal function, or other partial computation yielded no value. It returns a boolean and does not propagate $None$ as other operations do.

\subsection{Function Evaluation}

\begin{frameddef}[Function Evaluation]
Let $\mathbb{U}_{function} = \{ count, sum, min, max, avg, stddev, olag, olead, isnone \}$. The evaluation $eval_{function} : \mathbb{U}_{OCEL} \times \mathbb{U}_B \times \mathbb{U}_{function} \times ((\mathbb{U}_{lExpr})^* \cup \mathbb{U}_{opql}) \nrightarrow \mathbb{U}_V$ is:
\[
eval_{function}(L, b, f, p) = \begin{cases}
f(eval_{opql}(L, [b], p)) & \text{if } p \in \mathbb{U}_{opql} \\
f(L, (eval_{lExpr}(L, b, p_1), \dots, eval_{lExpr}(L, b, p_n))) & \text{if } p \in (\mathbb{U}_{lExpr})^*
\end{cases}
\]
\end{frameddef}

where $n = |p|$ in the second case. When a subquery is passed as argument, it is evaluated starting with the singleton table $[b]$ containing the current record, so the subquery may reference any name currently in scope. When expression arguments are passed, all are evaluated against the current record before the function is applied.

\section{Complete Grammar Reference}\label{sec:grammar-ref}

This section presents the complete OPQL grammar in EBNF. Terminals are shown in blue typewriter font via \term{...}; non-terminals in \textit{italic} via \nterm{...}. Concatenation is implicit (juxtaposition), alternatives use $|$, optionals use $[ \dots ]$, repetitions use $\{ \dots \}$, and grouping uses $( \dots )$.

\begin{framedgram}[Query Structure]
\begin{tabularx}{\linewidth}{l c X}
\nterm{entryPoint}     & = & \nterm{fullQuery} \\
\nterm{fullQuery}      & = & \{ \nterm{contextRule} \} \nterm{returnRule} \\
\nterm{contextRule}    & = & \nterm{patternRule} | \nterm{filterRule} | \nterm{keepRule} | \nterm{whenRule} \\
\end{tabularx}
\end{framedgram}

\begin{framedgram}[Pattern Syntax]
\begin{tabularx}{\linewidth}{l c X}
\nterm{patternRule}       & = & \term{PATTERN} \nterm{graphPatternList} [ \nterm{propositionalRule} ] \\
\nterm{graphPatternList}  & = & \nterm{graph} \{ \term{,} \nterm{graph} \} \\
\nterm{graph}             & = & \nterm{event} [ (\nterm{relationAny} | \nterm{relationRd}) \nterm{graphWithoutEvent} ] \\
                          & | & \nterm{object} [ (\nterm{relationAny} | \nterm{relationLd}) \nterm{graph} | \nterm{relationRd} \nterm{graphWithoutEvent} ] \\
\nterm{graphWithoutEvent} & = & \nterm{object} [ (\nterm{relationAny} | \nterm{relationLd}) \nterm{graph} ] \\
\nterm{event}             & = & \term{E(} [ \nterm{tag} ] [ \term{:} \nterm{name} ] \term{)} \\
\nterm{object}            & = & \term{O(} [ \nterm{tag} ] [ \term{:} \nterm{name} ] \term{)} \\
\nterm{relationAny}       & = & \term{-[} [ \nterm{tag} ] [ \term{:} \nterm{name} ] \term{]-} \\
\nterm{relationRd}        & = & \term{-[} [ \nterm{tag} ] [ \term{:} \nterm{name} ] \term{]->} \\
\nterm{relationLd}        & = & \term{<-[} [ \nterm{tag} ] [ \term{:} \nterm{name} ] \term{]-} \\
\nterm{tag}               & = & \nterm{symbolicName} \\
\nterm{name}              & = & \nterm{string} \\
\end{tabularx}
\end{framedgram}

\begin{framedgram}[Clauses]
\begin{tabularx}{\linewidth}{l c X}
\nterm{filterRule}        & = & \term{FILTER} \nterm{symbolicName} \{ \term{,} \nterm{symbolicName} \} \\
\nterm{keepRule}          & = & \term{KEEP} \nterm{projection} [ \nterm{propositionalRule} ] \\
\nterm{returnRule}        & = & \term{RETURN} \term{OCEL} \\
                          & | & \term{RETURN} \nterm{projection} [ \nterm{propositionalRule} ] \\
\nterm{whenRule}          & = & \term{WHEN} \nterm{expression} \term{AS} \nterm{symbolicName} \\
                          & & \quad [ \nterm{propositionalRule} ] \\
\nterm{propositionalRule} & = & (\term{SUBJECTTO} | \term{ST}) \nterm{expression} \\
\end{tabularx}
\end{framedgram}

\begin{framedgram}[Projection]
\begin{tabularx}{\linewidth}{l c X}
\nterm{projection}     & = & [ \term{DISTINCT} ] \term{*} [ \nterm{order} ] [ \nterm{limit} ] \\
                       & | & [ \term{DISTINCT} ] [ \term{*} \term{,} ] \nterm{projectionItem} \{ \term{,} \nterm{projectionItem} \} \\
                       & & \quad [ \nterm{order} ] [ \nterm{limit} ] \\
\nterm{projectionItem} & = & \nterm{sname} [ \nterm{binning} ] \\
\nterm{sname}          & = & \nterm{expression} [ \term{AS} \nterm{symbolicName} ] \\
\end{tabularx}
\end{framedgram}

\begin{framedgram}[Order and Limit]
\begin{tabularx}{\linewidth}{l c X}
\nterm{order}     & = & \term{ORDERBY} \nterm{orderItem} \{ \term{,} \nterm{orderItem} \} \\
\nterm{orderItem} & = & \nterm{expression} [ (\term{ASC} | \term{ASCENDING} | \term{DESC} | \term{DESCENDING}) ] \\
\nterm{limit}     & = & \term{LIMIT} \nterm{int} \\
\end{tabularx}
\end{framedgram}

\begin{framedgram}[Binning]
\begin{tabularx}{\linewidth}{l c X}
\nterm{binning}         & = & \term{BINNED(} \nterm{binningInterval} \{ \term{,} \nterm{binningInterval} \} \term{)} \\
\nterm{binningInterval} & = & (\term{(} | \term{[}) \nterm{intervalLimit} \term{,} \nterm{intervalLimit} \\
                        & & \quad (\term{)} | \term{]}) \term{AS} \nterm{intervalTarget} \\
\nterm{intervalLimit}   & = & \nterm{int} | \nterm{float} | \term{-inf} | [ \term{+} ] \term{inf} \\
\nterm{intervalTarget}  & = & \nterm{int} | \nterm{float} | \nterm{string} \\
\end{tabularx}
\end{framedgram}

\begin{framedgram}[Expressions]
\begin{tabularx}{\linewidth}{l c X}
\nterm{expression}  & = & \term{+} \nterm{expression} | \term{-} \nterm{expression} \\
                    & | & \nterm{expression} \term{\textasciicircum} \nterm{expression} \\
                    & | & \nterm{expression} (\term{*} | \term{/} | \term{\%}) \nterm{expression} \\
                    & | & \nterm{expression} (\term{+} | \term{-}) \nterm{expression} \\
                    & | & \nterm{expression} \nterm{compareSign} \nterm{expression} \\
                    & | & \term{NOT} \nterm{expression} \\
                    & | & \nterm{expression} \term{AND} \nterm{expression} \\
                    & | & \nterm{expression} \term{XOR} \nterm{expression} \\
                    & | & \nterm{expression} \term{OR} \nterm{expression} \\
                    & | & \term{(} \nterm{expression} \term{)} \\
                    & | & \nterm{valueType} \\
\nterm{compareSign} & = & \term{==} | \term{<=} | \term{>=} | \term{>} | \term{<} | \term{!=} \\
\end{tabularx}
\end{framedgram}

\begin{framedgram}[Values]
\begin{tabularx}{\linewidth}{l c X}
\nterm{valueType}         & = & \nterm{eoProperty} | \nterm{int} | \nterm{float} | \nterm{constantBool} | \term{None} \\
                          & | & \nterm{string} | \nterm{timestamp} | \nterm{duration} \\
                          & | & \nterm{symbolicName} | \nterm{rvFunctionCall} | \nterm{subquery} \\
\nterm{subquery}          & = & [ \term{MATERIALIZED} ] \term{(} \nterm{fullQuery} \term{)} \\
                          & | & \term{NOT MATERIALIZED} \term{(} \nterm{fullQuery} \term{)} \\
\nterm{eoProperty}        & = & \nterm{symbolicName} \term{[} \nterm{string} [ \term{@} \nterm{propertyTimestamp} ] \term{]} \\
\nterm{propertyTimestamp} & = & \nterm{symbolicName} | \nterm{timestamp} \\
\nterm{timestamp}         & = & \term{T(} \nterm{string} \term{)} \\
\nterm{duration}          & = & \term{D(} \nterm{int} \term{,} \nterm{int} \term{,} \nterm{int} \term{,} (\nterm{int} | \nterm{float}) \term{)} \\
\nterm{constantBool}      & = & \term{True} | \term{False} \\
\end{tabularx}
\end{framedgram}

\begin{framedgram}[Functions]
\begin{tabularx}{\linewidth}{l c X}
\nterm{rvFunctionCall} & = & \nterm{symbolicName}\term{(} \nterm{rvFunctionArg} \{ \term{,} \nterm{rvFunctionArg} \} \term{)} \\
\nterm{rvFunctionArg}  & = & \nterm{expression} | \nterm{fullQuery} \\
\end{tabularx}
\end{framedgram}

\begin{framedgram}[Lexical Elements]
\begin{tabularx}{\linewidth}{l c X}
\nterm{symbolicName} & = & (\term{a}..\term{z} | \term{A}..\term{Z}) \{ \term{a}..\term{z} | \term{A}..\term{Z} | \term{0}..\term{9} | \term{\_} \} \\
\nterm{string}       & = & \term{"} \{ \text{any encodable character} | \text{escape sequence} \} \term{"} \\
\nterm{int}          & = & \term{0} | (\term{1}..\term{9}) \{ \term{0}..\term{9} \} \\
\nterm{float}        & = & \nterm{int} \term{.} (\term{0}..\term{9}) \{ \term{0}..\term{9} \} \\
\end{tabularx}
\end{framedgram}

\section{Order and Limit}

Ordering and limiting are projection modifiers available in KEEP and RETURN clauses.

\begin{framedgram}[Order and Limit]
\begin{tabularx}{\linewidth}{l c X}
\nterm{order} & = & \opqltoken{ORDERBY} \nterm{orderItem} $\{ \term{,} \nterm{orderItem} \}$ \\
\nterm{orderItem} & = & \nterm{expression} [ (\opqltoken{ASC} | \opqltoken{ASCENDING} | \opqltoken{DESC} | \opqltoken{DESCENDING}) ] \\
\nterm{limit} & = & \opqltoken{LIMIT} \nterm{int} \\
\end{tabularx}
\end{framedgram}

An \nterm{order} clause specifies one or more sort keys. Each \nterm{orderItem} consists of an expression and an optional direction (ascending or descending). If no direction is specified, ascending order is the default. Multiple order items are applied in sequence: records that are indistinguishable under the first sort key are ordered by the second, and so on. If records remain indistinguishable after all sort keys, their relative order is implementation-defined.

A \nterm{limit} clause restricts the number of records in the result to the specified integer. Ordering is applied before limiting.

Both ordering and limiting operate on the table after all expressions in the projection have been evaluated and all filtering (via SUBJECTTO) has been applied.


\end{document}
